% Generated draft: requires mathematical and source review.
\documentclass[11pt]{article}
\usepackage[margin=25mm]{geometry}
\usepackage{fontspec}
\setmainfont{Latin Modern Roman}
\usepackage{amsmath,amssymb,mathtools,unicode-math}
\setmathfont{Latin Modern Math}
\usepackage{enumitem,needspace,xurl,hyperref}
\setlength{\parindent}{0pt}
\setlength{\parskip}{5pt}
\setlength{\emergencystretch}{3em}
\newcommand{\N}{\mathbb{N}}
\newcommand{\Z}{\mathbb{Z}}
\newcommand{\Q}{\mathbb{Q}}
\newcommand{\R}{\mathbb{R}}
\newcommand{\C}{\mathbb{C}}
\newcommand{\F}{\mathbb{F}}
\newcommand{\A}{\mathbb{A}}
\newcommand{\PP}{\mathbb{P}}
\newcommand{\E}{\mathbb{E}}
\newcommand{\eps}{\varepsilon}
\newcommand{\dd}{\,\mathrm{d}}
\begin{document}
\section*{UnsolvedMath ID: 30001117}
\textbf{Source problem code:} catalogue entry
\subsection{Definitions and mathematical statement}
\textbf{Source-level formulation.} For the shallow-slope martensitic-interface depinning model referred to in the supplied record, characterize analytically the different pinning regimes produced by the random nature of obstacles and their associated velocity exponents.

\textbf{Definitions and unresolved qualifications.} The supplied record does not specify the interface evolution equation, the spatial domain or dimension, initial or boundary conditions, the obstacle geometry or probability law, the driving parameter, or the admissible parameter ranges. It also does not define a pinning regime, the interface velocity, or the asymptotic relation that defines a velocity exponent. Consequently, the model name alone does not determine mathematical objects, hypotheses, or quantifiers sufficient for a precise classification or exponent statement. The formulation above preserves the supplied research question; a fully quantified mathematical statement requires the missing model specifications and definitions.

The supplied background additionally describes a general transition between regimes and qualifies its analytical understanding as model-dependent. The record does not specify which transitions, obstacle distributions, or parameter limits are intended, nor whether conclusions should hold almost surely, in expectation, or in another probabilistic sense. These choices remain unresolved.
\subsection{Short English statement}
Determine analytically how random obstacles produce different pinning regimes and associated velocity exponents for the shallow-slope martensitic-interface model. The supplied record omits the model equations and definitions needed to turn this question into a precise mathematical statement. Its partially solved label is supplied dated curation, not an independently verified conclusion.
\subsection{Sources}
\begin{itemize}
\item Local UnsolvedMath registry: supplied complete record for permanent ID \(30001117\), entitled \emph{Analytical Regimes of Martensitic-Interface Depinning}, in the Oberwolfach Reports Open Problems collection; exact source problem code \texttt{catalogue entry}. The record gives proposed year 2008 and names no individual proposer.
\item Supplied report citation: \emph{Interplay of Analysis and Probability in Physics} (2009), Oberwolfach Reports, pp.\ 3065--3138, DOI \texttt{10.4171/owr/2008/55}. Supplied literature-source label: ``original OWR report.'' \url{https://doi.org/10.4171/owr/2008/55}. The supplied background describes Oberwolfach Reports as documenting mathematical workshops and their open-problem discussions. No problem-specific page or passage is supplied.
\item Additional supplied literature source, labeled ``random pinning with correlated disorder'': \url{https://doi.org/10.1016/j.spa.2012.06.007}. The record supplies no authors, full bibliographic title, or precise result connecting this reference to the stated martensitic-interface model.
\item Supplied literature assessment, checked \texttt{2026-08-21}: status ``partially solved.'' The assessment states that rigorous pinning and depinning theories describe several obstacle and disorder regimes, while the general transition between the regimes posed in the report remains model-dependent and unresolved. This is supplied dated curation derived from the OWR curation, not an independently verified literature assessment; time-sensitive status and cited claims require independent verification.
\item The registry marks the statement as ``exact'' and reports that its mathematical content was checked against the official Oberwolfach report, with notation, context, and question form normalized without changing the source-level claim. These are supplied verification claims; no independent source consultation or verification is asserted here.
\end{itemize}
\end{document}
